\documentclass[11pt]{article}

\usepackage[a4paper,margin=1in]{geometry}
\usepackage{amsmath,amssymb,amsthm,mathtools}
\usepackage[final]{microtype}
\usepackage{braket}
\usepackage{enumitem}
\usepackage{hyperref}
\hypersetup{hidelinks}
\allowdisplaybreaks[2]

% ------------------------------------------------------------
% Macros
% ------------------------------------------------------------
\newcommand{\F}{\mathbb F}
\newcommand{\C}{\mathbb C}
\newcommand{\Z}{\mathbb Z}
\newcommand{\R}{\mathbb R}
\newcommand{\wt}{\operatorname{wt}}
\newcommand{\supp}{\operatorname{supp}}
\newcommand{\Span}{\operatorname{span}}
\newcommand{\diag}{\operatorname{diag}}
\newcommand{\Tr}{\operatorname{Tr}}
\newcommand{\id}{\mathbb I}
\newcommand{\KL}{\mathrm{KL}}
\newcommand{\BD}{\operatorname{BD}}
\newcommand{\pauli}{\mathcal P}
\newcommand{\ov}[1]{\overline{#1}}
\newcommand{\ketphi}{\ket{\phi}}
\newcommand{\ketpsi}{\ket{\psi}}
\newcommand{\ketzero}{\ket{0}}
\newcommand{\ketone}{\ket{1}}
\newcommand{\braidx}[1]{\bra{#1}}
\newcommand{\ketidx}[1]{\ket{#1}}

\newtheorem{theorem}{Theorem}[section]
\newtheorem{proposition}[theorem]{Proposition}
\newtheorem{lemma}[theorem]{Lemma}
\newtheorem{corollary}[theorem]{Corollary}
\newtheorem{definition}[theorem]{Definition}
\newtheorem{remark}[theorem]{Remark}

\title{Note on SS--LP Codes, Carry--Tag Lifts,\\
       and the Nonadditive Dyadic Bell Ladder}
\author{Internal working note}
\date{\today}

\begin{document}
\maketitle

\begin{abstract}
This note presents the theory behind the dyadic lift construction for two-dimensional
subset-sum/linear-programming (SS--LP) quantum codes.  The SS--LP framework fixes a
subset-sum residue structure so that a tensor product diagonal gate has a prescribed logical
diagonal action; the remaining problem is to choose amplitudes and supports satisfying the
Knill--Laflamme equations.  A carry-tag lift refines residues modulo $m$ to residues modulo
$2m$ by appending a tag system that records and corrects the dyadic carry bit.  This residue
arithmetic doubles the induced diagonal order, but distance preservation is a separate
certificate.

The main theorem treated here is the nonadditive dyadic Bell lift.  The lift appends two
qubits with weights $(m,m)$ and uses the Bell parity tags
$(\ket{00}+\ket{11})/\sqrt2$ and $(\ket{01}+\ket{10})/\sqrt2$.  Under a strong-sum
hypothesis, two unique anchor hypotheses, and a finite Knill--Laflamme check on the seed and
first lift, the Bell lift can be iterated indefinitely while preserving distance at least $3$.
The proof is given in detail.  Finally, the construction is interpreted as a residue-compatible
generalized concatenated quantum code (GCQC), and its relationship with stabilizer divisible-code
towers, CSS diagonal-gate synthesis, permutation-invariant/spin-code constructions, and recent
nonadditive diagonal-transversal searches, and fixed-distance scaling comparisons is summarized.
\end{abstract}

\begin{center}
\fbox{\parbox{0.92\textwidth}{\small
\textbf{Working note; a publication version will replace it.}  This is the
development note for the dyadic Bell ladder, included here so that the
repository is self-contained while the manuscript is being finished.  It states
and proves the ladder theorem in the abstract, and exhibits no concrete code:
the explicit \(((7,2,3))\) seed that discharges the theorem's hypotheses, the
resulting unconditional family, and the proof that its distance is exactly three
live in the accompanying code and Lean development, not in this file.  See the
repository \texttt{README.md} for what is verified and how.}}
\end{center}

\tableofcontents

% ============================================================
\section{Conventions and SS--LP background}
\label{sec:sslp-background}
% ============================================================

\subsection{Codes, distance, and the Knill--Laflamme convention}

All qubits are written in the computational basis.  For a bitstring
$x=(x_1,\ldots,x_n)\in\F_2^n$ we write $\ket{x}=\ket{x_1}\otimes\cdots\otimes\ket{x_n}$,
and $\wt(x)$ denotes Hamming weight.  A Pauli string has weight equal to the number of
non-identity tensor factors.  A two-dimensional code has logical basis
\[
  \ket{0_L},\qquad \ket{1_L}.
\]
It has distance at least $d$ in the error-detection sense if every Pauli error $E$ of weight
$<d$ compresses to a scalar on the code subspace~\cite{KnillLaflamme1997}:
\begin{equation}
  \bra{i_L}E\ket{j_L}=\lambda_E\delta_{ij},
  \qquad i,j\in\{0,1\}.
  \label{eq:kl-basic}
\end{equation}
For distance at least $3$, this means that \eqref{eq:kl-basic} is required for all Pauli strings
of weight at most $2$.  This is the convention used throughout the note.

\subsection{Subset-sum sectors and transversal diagonal phases}

The SS--LP construction begins by fixing a modulus $m\ge2$ and an integer angle vector~\cite{ZhangWuHuangZeng2025}
\[
  a=(a_1,\ldots,a_n),\qquad 0\le a_j\le m-1.
\]
For each bitstring $x\in\F_2^n$, define the integer subset sum and its residue modulo $m$ by
\begin{equation}
  A_a(x):=\sum_{j=1}^n a_jx_j,
  \qquad
  A_a(x)\bmod m\in\Z/m\Z .
  \label{eq:subset-sum}
\end{equation}
For a residue $b\in\Z/m\Z$, set
\begin{equation}
  S_b(a;m):=\{x\in\F_2^n: A_a(x)\equiv b\pmod m\}.
  \label{eq:subset-sum-block}
\end{equation}
The two-sector convention in this note is
\begin{equation}
  b_0=0,
  \qquad
  b_1=m-1.
  \label{eq:k2-residues}
\end{equation}
An SS--LP code chooses normalized logical states supported in these sectors:
\begin{equation}
  \ket{0_L}=\sum_{x\in S_0(a;m)} c^{(0)}_x\ket{x},
  \qquad
  \ket{1_L}=\sum_{x\in S_{m-1}(a;m)} c^{(1)}_x\ket{x}.
  \label{eq:sslp-codewords}
\end{equation}
The displayed sums allow coefficients to be zero; the actual support may be a proper subset of
the corresponding residue block.  Since the residue blocks are disjoint, the two logical states
are automatically orthogonal.

Let $\omega=e^{2\pi i/m}$ and define the physical transversal diagonal
\begin{equation}
  U_m(a):=\bigotimes_{j=1}^n
  \left(\ket{0}\!\bra{0}+\omega^{a_j}\ket{1}\!\bra{1}\right).
  \label{eq:sslp-physical-diagonal}
\end{equation}
A basis state in $S_b(a;m)$ is an eigenvector of $U_m(a)$ with eigenvalue $\omega^b$.  Therefore
\begin{equation}
  U_m(a)\ket{0_L}=\ket{0_L},
  \qquad
  U_m(a)\ket{1_L}=\omega^{-1}\ket{1_L}.
  \label{eq:sslp-logical-diagonal}
\end{equation}
The induced projective logical diagonal is
\begin{equation}
  D_m:=\diag(1,e^{-2\pi i/m}),
  \label{eq:Dm-def}
\end{equation}
which has projective order $m$.  The determinant-one representative
\begin{equation}
  R_m:=e^{\pi i/m}D_m=\diag(e^{\pi i/m},e^{-\pi i/m})
  \label{eq:Rm-def}
\end{equation}
has ordinary order $2m$.

The ``LP'' part of SS--LP refers to the fact that, once the subset-sum sectors have been fixed,
many Knill--Laflamme constraints become linear constraints on the probability weights
$|c_x^{(k)}|^2$, especially the constraints from diagonal $Z$-type errors.  In searches for
nonadditive codes with prescribed diagonal transversal gates, one first enumerates or designs the
integer data $(m,a,b_k)$, then solves the resulting marginal constraints and checks the remaining
Pauli constraints.  The residue support is responsible for the transversal logical phase, while the
KL equations are responsible for the distance.

\subsection{Optional self-complementarity and binary-dihedral symmetry}

Let $\ov{x}:=x\oplus 1^n$ be the bitwise complement.  A common additional symmetry is
\begin{equation}
  \ket{1_L}=X^{\otimes n}\ket{0_L}.
  \label{eq:self-complementarity}
\end{equation}
For this to be compatible with the residue convention \eqref{eq:k2-residues}, complementing a
string in $S_0(a;m)$ must send it to $S_{m-1}(a;m)$.  Since
\[
  A_a(\ov{x})\equiv \sum_{j=1}^n a_j-A_a(x)\pmod m,
\]
a sufficient and necessary residue condition is
\begin{equation}
  \sum_{j=1}^n a_j\equiv -1\pmod m.
  \label{eq:self-complementarity-residue}
\end{equation}
When \eqref{eq:self-complementarity} and \eqref{eq:self-complementarity-residue} hold, the
determinant-one diagonal $R_m$ together with the logical bit flip generates a binary-dihedral
logical group.  In this note the notation $\BD_{2m}$ is used for the group generated by the
order-$2m$ rotation $R_m$ and an order-four lift of the logical bit flip.  The dyadic distance
induction below does not require this symmetry; it is an optional enhancement of the cyclic
diagonal conclusion.

% ============================================================
\section{The carry--tag lift framework}
\label{sec:carry-tag-template}
% ============================================================

The purpose of a carry-tag lift is to refine an SS residue modulo $m$ to an SS residue modulo
$2m$.  The lift is best separated into two parts: an arithmetic phase-lift statement and an
independent distance certificate.

\subsection{Dyadic high bits}

Let $x\in S_{b_k}(a;m)$ for one of the logical residues $b_k$.  Reducing the integer subset sum
modulo $2m$, there is a unique high bit
\begin{equation}
  h_k(x)\in\{0,1\}
  \label{eq:high-bit-general}
\end{equation}
such that
\begin{equation}
  A_a(x)\equiv b_k+h_k(x)m\pmod{2m}.
  \label{eq:high-bit-residue}
\end{equation}
This bit is the dyadic carry that must be corrected by the appended tag.

\subsection{Two-state parity tags}

Let the tag consist of $t$ new qubits.  Choose a parity vector
\[
  \gamma=(\gamma_1,\ldots,\gamma_t)\in\F_2^t
\]
and append physical subset-sum weights
\[
  m\gamma_1,\ldots,m\gamma_t
\]
modulo $2m$.  A two-state parity tag is a pair of normalized states
\[
  \ket{\tau_0},\ket{\tau_1}\in(\C^2)^{\otimes t}
\]
whose supports lie in definite $\gamma$-parity sectors:
\begin{equation}
  \supp(\tau_p)\subseteq \{y\in\F_2^t: \gamma\cdot y\equiv p\pmod2\},
  \qquad p\in\{0,1\}.
  \label{eq:tag-parity-support}
\end{equation}
Choose bits $\delta_k\in\{0,1\}$, one for each logical sector.  The tagged lift is
\begin{equation}
  \ket{k_L^{\uparrow}}
  :=\sum_{x\in S_{b_k}(a;m)} c^{(k)}_x\ket{x}\ket{\tau_{h_k(x)\oplus\delta_k}}.
  \label{eq:two-state-tag-lift}
\end{equation}
The new modulus and angle vector are
\begin{equation}
  m^{\uparrow}=2m,
  \qquad
  a^{\uparrow}=(a_1,\ldots,a_n,m\gamma_1,\ldots,m\gamma_t).
  \label{eq:general-tag-angle}
\end{equation}
Every lifted basis string in the $k$th logical state has residue
\begin{align}
  A_{a^{\uparrow}}(x,y)
  &\equiv b_k+h_k(x)m+m\gamma\cdot y\pmod{2m} \nonumber\\
  &\equiv b_k+h_k(x)m+m(h_k(x)\oplus\delta_k)\pmod{2m} \nonumber\\
  &\equiv b_k+\delta_k m\pmod{2m}.
  \label{eq:tag-residue-computation}
\end{align}
For the two-sector convention $(b_0,b_1)=(0,m-1)$, the order-doubling choice is
\begin{equation}
  \delta_0=0,
  \qquad
  \delta_1=1.
  \label{eq:delta-order-doubling}
\end{equation}
Then the new residues are
\begin{equation}
  b_0^{\uparrow}=0,
  \qquad
  b_1^{\uparrow}=2m-1.
  \label{eq:new-residues-two-sector}
\end{equation}
Thus the induced logical diagonal changes from $D_m$ to $D_{2m}$.

\begin{theorem}[Phase-lift template]
\label{thm:phase-lift-template}
If a tagged lift places the lifted logical states in refined SS residue classes
$b_k^{\uparrow}\pmod {2m}$, then the physical diagonal $U_{2m}(a^{\uparrow})$ acts on the lifted
logical basis by the phases $e^{2\pi i b_k^{\uparrow}/(2m)}$.  In particular, in the two-sector
case \eqref{eq:new-residues-two-sector}, the induced projective logical diagonal is $D_{2m}$.
\end{theorem}

\begin{proof}
The proof is just the eigenvalue calculation.  For any computational basis vector
$\ket{z}$, the operator $U_{2m}(a^{\uparrow})$ gives the phase
$e^{2\pi i A_{a^{\uparrow}}(z)/(2m)}$.  If every basis vector in a logical state has the same
residue $b_k^{\uparrow}$ modulo $2m$, the whole logical state is an eigenvector with eigenvalue
$e^{2\pi i b_k^{\uparrow}/(2m)}$.  For residues $(0,2m-1)$ this is
$\diag(1,e^{-2\pi i/(2m)})=D_{2m}$.
\end{proof}

\subsection{Distance as a separate certificate}

The residue computation does not prove any distance statement.  To see the issue, decompose the
seed codewords into carry pieces
\begin{equation}
  \ket{k_L}=\ket{\psi_{k,0}}+\ket{\psi_{k,1}},
  \qquad
  \ket{\psi_{k,p}}:=\sum_{x:h_k(x)=p} c_x^{(k)}\ket{x}.
  \label{eq:carry-decomposition}
\end{equation}
Even if the unflagged seed satisfies KL, a lifted old-block Pauli $A\otimes\id$ sees only the
pairs of carry pieces that carry the same tag.  By \eqref{eq:two-state-tag-lift} the piece
$\ket{\psi_{k,p}}$ is tagged by $\tau_{p\oplus\delta_k}$, so pairing $\ket{\psi_{i,p}}$ with
$\ket{\psi_{j,q}}$ requires $p\oplus\delta_i=q\oplus\delta_j$, that is
$q=p\oplus\delta_i\oplus\delta_j$:
\begin{equation}
  \bra{i_L^{\uparrow}}A\otimes\id\ket{j_L^{\uparrow}}
  =\sum_p \bra{\psi_{i,p}}A\ket{\psi_{j,p\oplus\delta_i\oplus\delta_j}}
  \label{eq:old-pauli-sees-same-carry}
\end{equation}
when the tag states are orthogonal.  Note the shift $\delta_i\oplus\delta_j$: for the
order-doubling choice \eqref{eq:delta-order-doubling} it vanishes on the diagonal $i=j$ but
equals $1$ off it, so a lifted Pauli reads the seed's \emph{cross}-carry matrix elements, not
its same-carry ones.  This need not equal the original unflagged matrix element.
Distance preservation therefore requires additional information about how the seed KL data are
distributed across carry sectors and how low-weight tag errors act on the tag states.

A useful strong sufficient criterion is the following.

\begin{definition}[Carry-dephased KL condition]
\label{def:carry-dephased-kl}
Let $e=d-1$.  A carry decomposition \eqref{eq:carry-decomposition} is carry-dephased KL up to
weight $e$ relative to the tag offsets $\delta_0,\delta_1$ if, for every old-block Pauli $A$
with $\wt(A)\le e$,
\begin{equation}
  \sum_{p=0}^1 \bra{\psi_{i,p}}A\ket{\psi_{j,p\oplus\delta_i\oplus\delta_j}}=\lambda_A\delta_{ij}
  \label{eq:carry-dephased-kl}
\end{equation}
for a scalar $\lambda_A$ independent of the logical label.  The diagonal terms $i=j$ are the
same-carry sums $\sum_p\bra{\psi_{i,p}}A\ket{\psi_{i,p}}$; the off-diagonal ones are shifted.
\end{definition}

\begin{definition}[Tag invisibility]
\label{def:tag-invisibility}
A tag pair $\tau_0,\tau_1$ is invisible to order $e$ if, for every tag Pauli $B$ with
$\wt(B)\le e$,
\begin{equation}
  \bra{\tau_p}B\ket{\tau_q}=\mu_B\delta_{pq}
  \label{eq:tag-invisibility}
\end{equation}
for a scalar $\mu_B$ independent of $p$ and $q$.
\end{definition}

\begin{proposition}[Strong carry-tag lift criterion]
\label{prop:strong-carry-tag}
Fix $d\ge2$ and put $e=d-1$.  Suppose the seed carry decomposition is carry-dephased KL up to
weight $e$, and suppose the parity-compatible tag pair is invisible to order $e$.  Then the lifted
code \eqref{eq:two-state-tag-lift} satisfies the KL equations for every Pauli string of total
weight at most $e$.
\end{proposition}

\begin{proof}
Let $E=A\otimes B$ have total weight at most $e$, with $A$ acting on the old block and $B$ acting
on the tag block.  Expanding the lifted matrix element gives
\[
  \bra{i_L^{\uparrow}}E\ket{j_L^{\uparrow}}
  =\sum_{p,q}\bra{\psi_{i,p}}A\ket{\psi_{j,q}}
  \bra{\tau_{p\oplus\delta_i}}B\ket{\tau_{q\oplus\delta_j}}.
\]
By tag invisibility the tag factor is $\mu_B$ when $p\oplus\delta_i=q\oplus\delta_j$ and zero
otherwise, so the inner sum collapses to $q=p\oplus\delta_i\oplus\delta_j$.  Hence
\[
  \bra{i_L^{\uparrow}}E\ket{j_L^{\uparrow}}
  =\mu_B\sum_p \bra{\psi_{i,p}}A\ket{\psi_{j,p\oplus\delta_i\oplus\delta_j}}
  =\mu_B\lambda_A\delta_{ij},
\]
where the last equality is carry-dephased KL.  By linearity the same conclusion holds for
arbitrary linear combinations of such Pauli strings.
\end{proof}

The Bell lift below does not satisfy this strong criterion for $d=3$ because the two-qubit
operator $ZZ$ reads the Bell parity.  The point of the two-anchor Bell theorem is that distance
$3$ can still be propagated by a more refined cancellation mechanism.

% ============================================================
\section{The nonadditive dyadic Bell lift}
\label{sec:dyadic-bell-lift}
% ============================================================

\subsection{Definition of the lift}

The dyadic Bell lift is the two-qubit parity-tag lift with
\begin{equation}
  \gamma=(1,1),
  \qquad
  a^{\uparrow}=(a_1,\ldots,a_n,m,m),
  \qquad
  m^{\uparrow}=2m.
  \label{eq:bell-angle-update}
\end{equation}
Define
\begin{equation}
  \ketphi:=\frac{\ket{00}+\ket{11}}{\sqrt2},
  \qquad
  \ketpsi:=\frac{\ket{01}+\ket{10}}{\sqrt2},
  \label{eq:bell-states}
\end{equation}
and, equivalently,
\begin{equation}
  \ket{\Phi_p}:=\frac{1}{\sqrt2}
  \sum_{\substack{u,v\in\{0,1\}\\ u\oplus v=p}}\ket{uv},
  \qquad p\in\{0,1\}.
  \label{eq:bell-parity-states}
\end{equation}
Thus $\ket{\Phi_0}=\ketphi$ and $\ket{\Phi_1}=\ketpsi$.

For a seed of the form \eqref{eq:sslp-codewords}, define the high bit $h(x)$ by
\[
  A_a(x)\equiv b_k+h(x)m\pmod{2m},
  \qquad x\in S_{b_k}(a;m).
\]
The Bell lift is
\begin{equation}
  \boxed{
  \ket{k_L^{\uparrow}}
  =\sum_{x\in S_{b_k}(a;m)} c_x^{(k)}\ket{x}\ket{\Phi_{h(x)\oplus k}},
  \qquad k\in\{0,1\}.
  }
  \label{eq:bell-lift-definition}
\end{equation}
Equivalently, if
\begin{equation}
  \ket{k_L^{(h)}}:=\sum_{\substack{x\in S_{b_k}(a;m)\\ h(x)=h}}c_x^{(k)}\ket{x},
  \qquad h\in\{0,1\},
  \label{eq:low-high-pieces}
\end{equation}
then
\begin{equation}
  \ket{0_L^{\uparrow}}=\ket{0_L^{(0)}}\ketphi+\ket{0_L^{(1)}}\ketpsi,
  \qquad
  \ket{1_L^{\uparrow}}=\ket{1_L^{(0)}}\ketpsi+\ket{1_L^{(1)}}\ketphi.
  \label{eq:bell-normal-form-general}
\end{equation}
By Theorem~\ref{thm:phase-lift-template}, the lifted code is supported in residues
$0$ and $2m-1$ modulo $2m$, and its induced projective diagonal is $D_{2m}$.

\subsection{Iteration and logical phase growth}

Starting from a seed $(m_0,n_0,a^{(0)})$, define recursively
\begin{equation}
  m_{r+1}:=2m_r,
  \qquad
  n_{r+1}:=n_r+2,
  \qquad
  a^{(r+1)}:=(a^{(r)},m_r,m_r)\in(\Z/m_{r+1}\Z)^{n_{r+1}}.
  \label{eq:iterated-update}
\end{equation}
After $r$ lifts,
\begin{equation}
  m_r=2^r m_0,
  \qquad
  n_r=n_0+2r,
  \qquad
  a^{(r)}=(a^{(0)},m_0,m_0,2m_0,2m_0,\ldots,2^{r-1}m_0,2^{r-1}m_0).
  \label{eq:iterated-angle-vector}
\end{equation}
All entries of $a^{(r)}$ are viewed modulo $m_r$.  If $\mathcal C_r$ denotes the rung-$r$ code
space, then residue arithmetic gives
\begin{equation}
  U_{m_r}(a^{(r)})\big|_{\mathcal C_r}=\diag(1,e^{-2\pi i/m_r})=D_{m_r}.
  \label{eq:rung-r-diagonal}
\end{equation}
Thus each lift doubles the projective diagonal order.  If the optional self-complementary symmetry
persists, the cyclic diagonal growth is accompanied by binary-dihedral growth
$\BD_{2m_r}$.

% ============================================================
\section{The two-anchor Bell theorem and detailed proof}
\label{sec:bell-proof}
% ============================================================

This section proves the distance part of the dyadic Bell lift.  The proof uses a finite check on
rungs $0$ and $1$, then propagates KL from rung $r$ to rung $r+1$ for all $r\ge1$.

\subsection{Bell-pair Pauli matrix elements}

\begin{lemma}[Bell-pair Pauli table]
\label{lem:bell-pauli}
Let $B\in\{I,X,Y,Z\}^{\otimes2}$.  The only nonzero matrix elements within
$\Span\{\ketphi,\ketpsi\}$ are
\begin{equation}
\left(\bra{\phi}B\ket{\phi},\;\bra{\psi}B\ket{\psi},\;\bra{\phi}B\ket{\psi}\right)=
\begin{cases}
(1,1,0), & B\in\{II,XX\},\\
(-1,1,0), & B=YY,\\
(1,-1,0), & B=ZZ,\\
(0,0,1), & B\in\{IX,XI\},\\
(0,0,-i), & B\in\{YZ,ZY\},\\
(0,0,0), & \text{otherwise.}
\end{cases}
\label{eq:bell-pauli-table}
\end{equation}
For Paulis not listed in the first five lines, all four overlaps
$\bra{\phi}B\ket{\phi}$, $\bra{\psi}B\ket{\psi}$,
$\bra{\phi}B\ket{\psi}$, and $\bra{\psi}B\ket{\phi}$ vanish.
\end{lemma}

\begin{proof}
This is a direct calculation from \eqref{eq:bell-states}.  For instance, the diagonal overlaps
are obtained by applying $B$ to the computational-basis terms in $\ket{00}+\ket{11}$ and
$\ket{01}+\ket{10}$ and retaining only the components that remain in the same two-dimensional
span.  The off-diagonal overlaps are computed similarly.  Exhausting the sixteen two-qubit
Paulis gives exactly the cases in \eqref{eq:bell-pauli-table}; Hermiticity gives the conjugate
off-diagonal entries.
\end{proof}

\subsection{Two-value subset sums and small-weight separation}

The theorem assumes a strong total-sum condition on the seed.  This condition ensures that, on
each rung, every residue block has only two possible integer subset-sum values.

\begin{lemma}[Two-value subset sums]
\label{lem:two-value-subsets}
Let $a_i\in\{0,1,\ldots,m-1\}$ and let $A_{\rm tot}:=\sum_i a_i$.  If
$A_{\rm tot}<2m$, then for any residue $b\in\{0,1,\ldots,m-1\}$ and any
$x\in S_b(a;m)$,
\begin{equation}
  A_a(x)\in\{b,b+m\}.
  \label{eq:two-values}
\end{equation}
Equivalently, each subset-sum block has a well-defined carry bit read directly from the integer
subset sum.
\end{lemma}

\begin{proof}
The integer sum satisfies $0\le A_a(x)\le A_{\rm tot}<2m$.  If it is congruent to $b$ modulo
$m$, the only possible values in the interval $[0,2m)$ are $b$ and $b+m$.
\end{proof}

The strong-sum identity used in the theorem is
\begin{equation}
  \sum_{j=1}^{n_0}a_j=2m_0-1.
  \label{eq:strong-sum-identity}
\end{equation}
Under the update \eqref{eq:iterated-update}, it propagates as
\begin{equation}
  \sum_{j=1}^{n_r}a^{(r)}_j=2m_r-1.
  \label{eq:strong-sum-propagates}
\end{equation}
Indeed,
\[
  (2m_r-1)+2m_r=2m_{r+1}-1.
\]
Hence Lemma~\ref{lem:two-value-subsets} applies at every rung.

\begin{lemma}[No small strings in the opposite residue after the first lift]
\label{lem:no-small-weight-Smr1}
Let $m_r=2^r m_0$ and let $a^{(r)}$ be given by \eqref{eq:iterated-angle-vector}.  Assume
\begin{equation}
  0\le a_j\le m_0-2
  \qquad\text{and}\qquad
  \sum_{j=1}^{n_0}a_j=2m_0-1.
  \label{eq:no-small-weight-hyp}
\end{equation}
Then for every $r\ge1$:
\begin{enumerate}[label=\textup{(\roman*)}]
\item no string of Hamming weight $1$ or $2$ lies in $S_{m_r-1}(a^{(r)};m_r)$;
\item equivalently, no string at Hamming distance at most $2$ from $1^{n_r}$ lies in
$S_0(a^{(r)};m_r)$.
\end{enumerate}
\end{lemma}

\begin{proof}
We first prove (i).  Let $x$ have Hamming weight $1$.  Its subset sum is a single component of
$a^{(r)}$.  Every component of $a^{(r)}$ is at most $m_r/2$: seed entries are at most
$m_0-2<m_r/2$, and appended entries are $m_0,2m_0,\ldots,2^{r-1}m_0=m_r/2$.  Since
$m_r-1>m_r/2$, no weight-one string has residue $m_r-1$.

Now let $x$ have Hamming weight $2$, and let its integer subset sum be $x_1+x_2$, where
$x_1,x_2$ are two components of $a^{(r)}$.  Again $0\le x_1,x_2\le m_r/2$, so
$0\le x_1+x_2\le m_r$.  If the residue is $m_r-1$, this range forces
\begin{equation}
  x_1+x_2=m_r-1.
  \label{eq:weight-two-eq}
\end{equation}
Reducing \eqref{eq:weight-two-eq} modulo $m_0$ gives $x_1+x_2\equiv -1\pmod{m_0}$.  Appended
components are multiples of $m_0$, while seed components lie in $\{0,1,\ldots,m_0-2\}$.  If
both $x_1,x_2$ are appended components, the congruence is impossible.  If exactly one is a seed
component, that seed component would have to be $m_0-1$, which is excluded.  If both are seed
components, then $x_1+x_2\le2m_0-4$, while $m_r-1\ge2m_0-1$ for $r\ge1$, again impossible.  This
proves (i).

For (ii), use \eqref{eq:strong-sum-propagates}.  For any bitstring $x$,
\[
  A_{a^{(r)}}(\ov{x})\equiv \sum_j a_j^{(r)}-A_{a^{(r)}}(x)
  \equiv -1-A_{a^{(r)}}(x)\pmod{m_r}.
\]
Thus $x\in S_{m_r-1}$ if and only if $\ov{x}\in S_0$.  Complementation maps weight at most $2$
to distance at most $2$ from $1^{n_r}$, proving (ii).
\end{proof}

\subsection{The theorem}

\begin{theorem}[Two-anchor Bell-flag ladder theorem]
\label{thm:two-anchor-bell}
Fix $m_0\ge2$ and $n_0\ge1$.  Let
$a^{(0)}=(a_1,\ldots,a_{n_0})$ have integer representatives satisfying
\begin{equation}
  0\le a_j\le m_0-2,
  \qquad
  \sum_{j=1}^{n_0}a_j=2m_0-1.
  \label{eq:theorem-angle-hypotheses}
\end{equation}
Let normalized seed states $\ket{0^{(0)}}$ and $\ket{1^{(0)}}$ be supported on
$S_0(a^{(0)};m_0)$ and $S_{m_0-1}(a^{(0)};m_0)$, respectively.  Assume:
\begin{enumerate}[label=\textup{(\roman*)}]
\item \textup{(unique bottom anchor)} within $\supp(\ket{0^{(0)}})$, the only basis string whose
integer subset sum is $0$ is $0^{n_0}$; equivalently, every other support string in
$\ket{0^{(0)}}$ has integer subset sum $m_0$;
\item \textup{(unique top anchor)} within $\supp(\ket{1^{(0)}})$, the only basis string whose
integer subset sum is $2m_0-1$ is $1^{n_0}$; equivalently, every other support string in
$\ket{1^{(0)}}$ has integer subset sum $m_0-1$;
\item the seed rung $r=0$ and the first Bell lift $r=1$ both satisfy the Knill--Laflamme
equations for all Pauli errors of weight at most $2$.
\end{enumerate}
Then every rung obtained by iterating the dyadic Bell lift satisfies the KL equations for all
Pauli errors of weight at most $2$.  In particular, every rung has distance at least $3$.
Moreover, the induced transversal diagonal on rung $r$ is $D_{m_r}$ with $m_r=2^r m_0$.  If the
seed is self-complementary in the sense of \eqref{eq:self-complementarity}, then every rung is
self-complementary and admits the corresponding binary-dihedral transversal group $\BD_{2m_r}$.
\end{theorem}

The proof is broken into structural lemmas followed by the KL induction.

\subsection{Bell recursion and anchor structure}

Choose phases of the two logical seed states so that the coefficients of the anchors
$0^{n_0}$ and $1^{n_0}$ are positive real.  Write
\begin{equation}
  \ket{0^{(r)}}=\alpha_r\ket{0^{n_r}}+\ket{u_r},
  \qquad
  \ket{1^{(r)}}=\beta_r\ket{1^{n_r}}+\ket{v_r},
  \label{eq:anchor-decomposition}
\end{equation}
where $\ket{u_r}$ and $\ket{v_r}$ contain no anchor component.

\begin{lemma}[Bell recursion with two anchors]
\label{lem:bell-recursion}
Under the hypotheses of Theorem~\ref{thm:two-anchor-bell}, for every $r\ge0$:
\begin{enumerate}[label=\textup{(\roman*)}]
\item the unique bottom and top anchors persist on rung $r$;
\item the lift from rung $r$ to rung $r+1$ has the Bell normal form
\begin{equation}
  \ket{0^{(r+1)}}=\alpha_r\ket{0^{n_r}}\ketphi+\ket{u_r}\ketpsi,
  \qquad
  \ket{1^{(r+1)}}=\beta_r\ket{1^{n_r}}\ketphi+\ket{v_r}\ketpsi;
  \label{eq:bell-recursion-r}
\end{equation}
\item the anchor coefficients satisfy
\begin{equation}
  \alpha_{r+1}=\frac{\alpha_r}{\sqrt2},
  \qquad
  \beta_{r+1}=\frac{\beta_r}{\sqrt2};
  \label{eq:anchor-halving}
\end{equation}
\item $\alpha_r=\beta_r$ for every $r\ge0$;
\item the non-anchor parts obey
\begin{equation}
  \ket{u_{r+1}}=\alpha_{r+1}\ket{0^{n_r}11}+\ket{u_r}\ketpsi,
  \qquad
  \ket{v_{r+1}}=\beta_{r+1}\ket{1^{n_r}00}+\ket{v_r}\ketpsi.
  \label{eq:uv-recursion}
\end{equation}
\end{enumerate}
\end{lemma}

\begin{proof}
The strong-sum identity propagates by \eqref{eq:strong-sum-propagates}, so Lemma~\ref{lem:two-value-subsets}
applies on every rung.  In the bottom sector, integer subset sums are $0$ or $m_r$; in the top
sector, they are $m_r-1$ or $2m_r-1$.

The unique anchors persist by induction.  The bottom anchor $0^{n_r}$ has carry bit $0$, so the
lift appends the even parity pair.  The branch $0^{n_r}00$ has new subset sum $0$, while the
branch $0^{n_r}11$ has new subset sum $2m_r=m_{r+1}$.  Every other bottom support string has old
subset sum $m_r$ and carry bit $1$, so the lift appends odd parity and the new total subset sum
is $m_{r+1}$.  Hence the only bottom support string with subset sum $0$ on rung $r+1$ is
$0^{n_{r+1}}$.

For the top sector, the top anchor $1^{n_r}$ has old subset sum $2m_r-1$ and carry bit $1$.  Since
the logical label is $k=1$, the Bell rule appends even parity.  The branch $1^{n_r}00$ has new
subset sum $m_{r+1}-1$, while $1^{n_r}11$ has new subset sum $2m_{r+1}-1$.  Every other top
support string has old subset sum $m_r-1$ and carry bit $0$, and the parity shift again appends
odd parity; the new total subset sum is $m_{r+1}-1$.  Hence the only top support string with
subset sum $2m_{r+1}-1$ is $1^{n_{r+1}}$.

The same reasoning gives the Bell normal form \eqref{eq:bell-recursion-r}: anchors receive
$\ketphi$ and non-anchors receive $\ketpsi$.  Expanding
$\ketphi=(\ket{00}+\ket{11})/\sqrt2$, the coefficients of the new anchors
$0^{n_r}00$ and $1^{n_r}11$ are $\alpha_r/\sqrt2$ and $\beta_r/\sqrt2$, proving
\eqref{eq:anchor-halving}.  Subtracting the new anchors from \eqref{eq:bell-recursion-r} gives
\eqref{eq:uv-recursion}.

It remains to prove equality of the anchor coefficients.  On the first lifted rung, let
$E_0=\id^{\otimes n_0}\otimes ZZ$ act on the appended Bell pair.  From
\eqref{eq:bell-recursion-r} with $r=0$ and Lemma~\ref{lem:bell-pauli},
\[
  \bra{0^{(1)}}E_0\ket{0^{(1)}}
  =\alpha_0^2\bra{\phi}ZZ\ket{\phi}
  +\braket{u_0|u_0}\bra{\psi}ZZ\ket{\psi}
  =2\alpha_0^2-1.
\]
Similarly,
\[
  \bra{1^{(1)}}E_0\ket{1^{(1)}}=2\beta_0^2-1.
\]
The first-lift KL hypothesis makes these two diagonal matrix elements equal, so
$\alpha_0^2=\beta_0^2$.  The phases were chosen so that both are positive real, hence
$\alpha_0=\beta_0$.  The halving relation \eqref{eq:anchor-halving} propagates this equality to
every rung.
\end{proof}

\begin{lemma}[Unique low-weight buddies]
\label{lem:buddy-unique}
Under the hypotheses of Theorem~\ref{thm:two-anchor-bell}, for every $r\ge1$:
\begin{enumerate}[label=\textup{(\roman*)}]
\item the only basis string in $\supp(\ket{u_r})$ with Hamming weight at most $2$ is
\begin{equation}
  q_r:=0^{n_r-2}11,
  \label{eq:qr-def}
\end{equation}
and its coefficient is $\alpha_r$;
\item the only basis string in $\supp(\ket{v_r})$ within Hamming distance at most $2$ from
$1^{n_r}$ is
\begin{equation}
  p_r:=1^{n_r-2}00,
  \label{eq:pr-def}
\end{equation}
and its coefficient is $\beta_r$.
\end{enumerate}
\end{lemma}

\begin{proof}
We prove both claims by induction.  For $r=1$, the recursion gives
\[
  \ket{u_1}=\alpha_1\ket{0^{n_0}11}+\ket{u_0}\ketpsi.
\]
Every basis string in $\supp(\ket{u_0})$ has Hamming weight at least $2$: a weight-one string in
the bottom residue block would have subset sum $0$ only if its single nonzero position had
$a_j=0$, and such strings are excluded from the support by the unique bottom-anchor condition.
Tensoring with $\ketpsi$ appends a pair of Hamming weight $1$, so every basis string in
$\supp(\ket{u_0}\ketpsi)$ has weight at least $3$.  Hence the unique weight-at-most-two string in
$\supp(\ket{u_1})$ is $0^{n_0}11$, with coefficient $\alpha_1$.

Similarly,
\[
  \ket{v_1}=\beta_1\ket{1^{n_0}00}+\ket{v_0}\ketpsi.
\]
Any basis string in $\supp(\ket{v_0})$ has distance at least $2$ from $1^{n_0}$.  A distance-one
string from $1^{n_0}$ would have subset sum $2m_0-1-a_j$; this cannot be $m_0-1$ unless
$a_j=m_0$, impossible, and cannot be $2m_0-1$ unless $a_j=0$, which is excluded from the support
by the unique top-anchor condition.  Tensoring with $\ketpsi$ appends a pair at distance $1$ from
$11$, so all those strings have distance at least $3$ from $1^{n_1}$.  Thus the only string in
$\supp(\ket{v_1})$ within distance at most $2$ from $1^{n_1}$ is $1^{n_0}00$, with coefficient
$\beta_1$.

For the induction step, use \eqref{eq:uv-recursion}.  If a string $t\in\supp(\ket{u_r})$ is
tensored with a basis string in $\ketpsi$, the Hamming weight increases by $1$.  By the induction
hypothesis, either $t=q_r$ has weight $2$, or $t$ has weight at least $3$.  Therefore every string
in $\supp(\ket{u_r}\ketpsi)$ has weight at least $3$, and the only weight-at-most-two string in
$\ket{u_{r+1}}$ is the newly created $0^{n_r}11$, with coefficient $\alpha_{r+1}$.  The argument
for $v_r$ is the complement: tensoring with $\ketpsi$ increases the distance from the all-one
string by $1$, so the only near-top string on rung $r+1$ is $1^{n_r}00$, with coefficient
$\beta_{r+1}$.
\end{proof}

\begin{lemma}[Anchor-leakage cancellation]
\label{lem:anchor-leakage}
Under the hypotheses of Theorem~\ref{thm:two-anchor-bell}, fix $r\ge1$ and let $A$ be an
$n_r$-qubit Pauli operator with $\wt(A)\le2$.  Then:
\begin{enumerate}[label=\textup{(\roman*)}]
\item $\bra{0^{n_r}}A\ket{u_r}=0$ unless $A$ acts nontrivially exactly on the last two qubits of
rung $r$, with $X$ or $Y$ on each of those two qubits.  In the exceptional case,
$A=\id^{\otimes(n_r-2)}\otimes P$ with $P\in\{XX,XY,YX,YY\}$, and
\begin{equation}
  \bra{0^{n_r}}A\ket{u_r}=\alpha_r\bra{00}P\ket{11};
  \label{eq:anchor-leakage-zero}
\end{equation}
\item $\bra{1^{n_r}}A\ket{v_r}=0$ unless the same exceptional form holds, and then
\begin{equation}
  \bra{1^{n_r}}A\ket{v_r}=\beta_r\bra{11}P\ket{00}
  =\beta_r\overline{\bra{00}P\ket{11}};
  \label{eq:anchor-leakage-one}
\end{equation}
\item the real parts agree:
\begin{equation}
  \Re\!\left(\alpha_r\bra{0^{n_r}}A\ket{u_r}\right)
  =
  \Re\!\left(\beta_r\bra{1^{n_r}}A\ket{v_r}\right).
  \label{eq:anchor-leakage-real}
\end{equation}
\end{enumerate}
\end{lemma}

\begin{proof}
If $\bra{0^{n_r}}A\ket{u_r}\ne0$, then $A$ maps some support string of $u_r$ to $0^{n_r}$.  Since
$\wt(A)\le2$, that support string must have Hamming weight at most $2$.  By
Lemma~\ref{lem:buddy-unique}, the only possibility is $q_r=0^{n_r-2}11$, with coefficient
$\alpha_r$.  Mapping $q_r$ to $0^{n_r}$ requires flipping the last two bits and acting trivially
elsewhere, so the last-two-qubit Pauli must belong to $\{XX,XY,YX,YY\}$.  This proves (i).

The proof of (ii) is identical using $p_r=1^{n_r-2}00$.  Mapping $p_r$ to $1^{n_r}$ requires
flipping the last two bits, and the matrix element is
$\beta_r\bra{11}P\ket{00}$.  Since $P$ is Hermitian, this is
$\beta_r\overline{\bra{00}P\ket{11}}$.  Finally, Lemma~\ref{lem:bell-recursion} gives
$\alpha_r=\beta_r$, so the real parts in (iii) agree.
\end{proof}

\subsection{Proof of the theorem}

\begin{proof}[Proof of Theorem~\ref{thm:two-anchor-bell}]
Rungs $0$ and $1$ satisfy the KL equations by hypothesis.  We prove by induction that if rung
$r\ge1$ satisfies KL for all Pauli errors of weight at most $2$, then rung $r+1$ also satisfies
KL for all such errors.

Let $E$ be a Pauli string on $n_{r+1}=n_r+2$ qubits with $\wt(E)\le2$.  Write
\begin{equation}
  E=A\otimes B,
  \label{eq:EAB}
\end{equation}
where $A$ acts on the old $n_r$ qubits and $B$ acts on the newly appended Bell pair.  Then
$\wt(A)\le2-\wt(B)$.  By Lemma~\ref{lem:bell-recursion},
\begin{equation}
  \ket{0^{(r+1)}}=\alpha_r\ket{0^{n_r}}\ketphi+\ket{u_r}\ketpsi,
  \qquad
  \ket{1^{(r+1)}}=\beta_r\ket{1^{n_r}}\ketphi+\ket{v_r}\ketpsi,
  \label{eq:proof-bell-recursion}
\end{equation}
with $\alpha_r=\beta_r$.

\medskip
\noindent\textbf{Case 1: $B$ has no matrix elements inside $\Span\{\ketphi,\ketpsi\}$.}
If
\[
  B\notin\{II,IX,XI,XX,YY,ZZ,YZ,ZY\},
\]
then Lemma~\ref{lem:bell-pauli} says that every Bell overlap appearing in
$\bra{i^{(r+1)}}A\otimes B\ket{j^{(r+1)}}$ is zero.  Hence all KL matrix elements are zero, so KL
holds with $\lambda_E=0$.

\medskip
\noindent\textbf{Case 2: $B=XX$.}
Here $\wt(B)=2$, so $A=\id$.  Since $XX$ acts as the identity on both $\ketphi$ and $\ketpsi$,
$\id\otimes XX$ acts as the identity on the lifted code space.  KL holds with $\lambda_E=1$.

\medskip
\noindent\textbf{Case 3: $B\in\{ZZ,YY\}$.}
Again $\wt(B)=2$, so $A=\id$.  The Bell table shows that $ZZ$ and $YY$ are diagonal in the
$\{\ketphi,\ketpsi\}$ basis.  Therefore the diagonal logical expectations depend only on the
anchor weight and the non-anchor norm.  Since $\alpha_r=\beta_r$ and both logical states are
normalized, the two diagonal expectations are equal.

For the off-diagonal matrix element, first note that $\braket{u_r|v_r}=0$.  Indeed,
orthogonality of the rung-$r$ logical states gives
\[
0=\braket{0^{(r)}|1^{(r)}}
=\alpha_r\beta_r\braket{0^{n_r}|1^{n_r}}
+\alpha_r\braket{0^{n_r}|v_r}
+\beta_r\braket{u_r|1^{n_r}}
+\braket{u_r|v_r}.
\]
The first term vanishes because the anchors are distinct, and the middle two vanish because the
bottom and top residue sectors are disjoint.  Hence $\braket{u_r|v_r}=0$.  Since $ZZ$ and $YY$
have no off-diagonal Bell overlaps, the off-diagonal KL matrix element is zero.

\medskip
\noindent\textbf{Case 4: $B\in\{YZ,ZY\}$.}
Here $\wt(B)=2$ and $A=\id$.  By Lemma~\ref{lem:bell-pauli}, these Paulis have only off-diagonal
Bell overlaps.  In diagonal logical matrix elements, the old-block overlaps that can appear are
$\braket{0^{n_r}|u_r}$, $\braket{u_r|0^{n_r}}$, $\braket{1^{n_r}|v_r}$, and
$\braket{v_r|1^{n_r}}$, all of which vanish by the definition of the non-anchor parts.  In the
off-diagonal logical matrix element, the possible old-block overlaps are
$\braket{0^{n_r}|v_r}$ and $\braket{u_r|1^{n_r}}$, which vanish by residue-sector disjointness.
Thus every KL matrix element in this case is zero.

\medskip
\noindent\textbf{Case 5: $B\in\{IX,XI\}$.}
Here $\wt(B)=1$, so $\wt(A)\le1$.  The Bell table says that $B$ swaps $\ketphi$ and $\ketpsi$ with
unit amplitude.  For the logical-$0$ diagonal element,
\begin{align*}
  \bra{0^{(r+1)}}A\otimes B\ket{0^{(r+1)}}
  &=\alpha_r\bra{0^{n_r}}A\ket{u_r}
    +\alpha_r\bra{u_r}A\ket{0^{n_r}} \\
  &=2\Re\!\left(\alpha_r\bra{0^{n_r}}A\ket{u_r}\right).
\end{align*}
Since $\wt(A)\le1$, Lemma~\ref{lem:anchor-leakage} gives
$\bra{0^{n_r}}A\ket{u_r}=0$, because the exceptional leakage requires flipping the last two
qubits.  The logical-$1$ diagonal element vanishes in the same way.

For the off-diagonal element,
\[
  \bra{0^{(r+1)}}A\otimes B\ket{1^{(r+1)}}
  =\alpha_r\bra{0^{n_r}}A\ket{v_r}
   +\beta_r\bra{u_r}A\ket{1^{n_r}}.
\]
Both terms vanish.  If the first were nonzero, a support string of $v_r$ would have to be within
Hamming distance at most $1$ of $0^{n_r}$ and lie in $S_{m_r-1}$, contradicting
Lemma~\ref{lem:no-small-weight-Smr1}.  The second term is ruled out by the complement form of the
same lemma.  Hence KL holds with scalar zero.

\medskip
\noindent\textbf{Case 6: $B=II$.}
Here $E=A\otimes\id$ acts only on the old block and $\wt(A)\le2$.

For the off-diagonal KL condition, orthogonality of $\ketphi$ and $\ketpsi$ gives
\begin{equation}
  \bra{0^{(r+1)}}A\otimes\id\ket{1^{(r+1)}}
  =\alpha_r\beta_r\bra{0^{n_r}}A\ket{1^{n_r}}
   +\bra{u_r}A\ket{v_r}.
  \label{eq:lifted-offdiag-oldblock}
\end{equation}
The rung-$r$ off-diagonal KL equation for $A$ says
\begin{align}
0&=\bra{0^{(r)}}A\ket{1^{(r)}} \nonumber\\
 &=\alpha_r\beta_r\bra{0^{n_r}}A\ket{1^{n_r}}
 +\alpha_r\bra{0^{n_r}}A\ket{v_r}
 +\beta_r\bra{u_r}A\ket{1^{n_r}}
 +\bra{u_r}A\ket{v_r}.
 \label{eq:rung-r-offdiag-expanded}
\end{align}
The two cross-anchor terms vanish by Lemma~\ref{lem:no-small-weight-Smr1}: a weight-at-most-two
Pauli cannot connect the bottom anchor to the top residue block, nor the bottom residue block to
the top anchor.  The remaining two terms are exactly
\eqref{eq:lifted-offdiag-oldblock}, so the lifted off-diagonal matrix element is zero.

For the diagonal KL condition, again using $\braket{\phi|\psi}=0$,
\begin{align}
  \bra{0^{(r+1)}}A\otimes\id\ket{0^{(r+1)}}
  &=\alpha_r^2\bra{0^{n_r}}A\ket{0^{n_r}}
    +\bra{u_r}A\ket{u_r},
    \label{eq:lifted-diag-zero}\\
  \bra{1^{(r+1)}}A\otimes\id\ket{1^{(r+1)}}
  &=\beta_r^2\bra{1^{n_r}}A\ket{1^{n_r}}
    +\bra{v_r}A\ket{v_r}.
    \label{eq:lifted-diag-one}
\end{align}
We must show the two expressions are equal.  The rung-$r$ diagonal KL equation gives
\begin{align}
0
&=\bra{0^{(r)}}A\ket{0^{(r)}}-\bra{1^{(r)}}A\ket{1^{(r)}} \nonumber\\
&=\alpha_r^2\bra{0^{n_r}}A\ket{0^{n_r}}
 -\beta_r^2\bra{1^{n_r}}A\ket{1^{n_r}}
 +\bra{u_r}A\ket{u_r}-\bra{v_r}A\ket{v_r} \nonumber\\
&\quad +2\Re\!\left(\alpha_r\bra{0^{n_r}}A\ket{u_r}
 -\beta_r\bra{1^{n_r}}A\ket{v_r}\right).
 \label{eq:rung-r-diag-expanded}
\end{align}
By Lemma~\ref{lem:anchor-leakage}, the real parts in the final line are equal, so the final line
of \eqref{eq:rung-r-diag-expanded} is zero.  What remains is exactly the difference between
\eqref{eq:lifted-diag-zero} and \eqref{eq:lifted-diag-one}.  Hence the two lifted diagonal matrix
elements are equal.

These six cases exhaust all possible restrictions of a weight-at-most-two Pauli to the newest
Bell block.  Therefore rung $r+1$ satisfies KL for all Pauli errors of weight at most $2$.  The
induction begins at $r=1$ by hypothesis, so every rung has distance at least $3$.

The diagonal logical action was proved in Section~\ref{sec:dyadic-bell-lift}.  Finally, if the
seed is self-complementary, the strong-sum identity implies that complementing a support string
flips the carry bit.  Since $X^{\otimes2}$ fixes both $\ketphi$ and $\ketpsi$, the relation
$\ket{1^{(r)}}=X^{\otimes n_r}\ket{0^{(r)}}$ propagates through the lift.  Combining this
transversal bit flip with the diagonal $D_{m_r}$ gives the stated binary-dihedral conclusion.
\end{proof}

% ============================================================
\section{Generalized concatenated quantum code viewpoint}
\label{sec:gcqc}
% ============================================================

\subsection{GCQC reminder}

Generalized concatenated quantum codes decompose an inner Hilbert space into mutually orthogonal
subcodes~\cite{GrasslShorSmithSmolinZeng2009}
\begin{equation}
  B=\bigoplus_{\mu\in\mathcal A}B_\mu,
  \label{eq:gcqc-inner-decomposition}
\end{equation}
and let the outer code select subcodes rather than single inner basis states.  In ordinary
concatenation, an outer alphabet symbol is mapped to a basis state or fixed codeword of the inner
code.  In generalized concatenation, an outer symbol may select an entire subspace $B_\mu$.
This additional freedom is useful for both stabilizer and nonadditive constructions.

At the level of a single carry-tag lift, the seed is decomposed into sectors indexed by a carry
label $\mu$, and the tag supplies mutually orthogonal tag subspaces $T_\mu$.  A general tagged
state has the form
\begin{equation}
  \ket{\Psi_i}=\sum_{\mu}\ket{\psi_{i,\mu}}\otimes\ket{\tau_{i,\mu}},
  \label{eq:general-tagged-state}
\end{equation}
where $\ket{\psi_{i,\mu}}$ is the old-block component and $\ket{\tau_{i,\mu}}$ lies in the tag
subspace selected by $\mu$ and possibly by the logical label $i$.

For a Pauli $A\otimes B$, the KL matrix element expands as
\begin{equation}
  \bra{\Psi_i}A\otimes B\ket{\Psi_j}
  =\sum_{\mu,\nu}\bra{\psi_{i,\mu}}A\ket{\psi_{j,\nu}}
  \bra{\tau_{i,\mu}}B\ket{\tau_{j,\nu}}.
  \label{eq:gcqc-kl-expansion}
\end{equation}
This is the GCQC form of the distance problem: distance depends on a compatibility between the
old sector-by-sector KL data and the low-weight error table of the tag subcodes.

\subsection{Carry-tag lifts as residue-compatible GCQC}

Carry-tag lifts are GCQC with an additional transversal-residue constraint.  The sector label is
not an arbitrary alphabet symbol; it is a subset-sum carry bit chosen so that the refined code lies
in prescribed residue classes modulo $2m$.  The tag subcodes are not arbitrary inner subcodes; they
must have support in parity sectors that contribute the right multiple of $m$ to the subset sum.
Thus a carry-tag lift has three logically separate layers:
\begin{equation}
\begin{gathered}
  \text{old-code sector decomposition}
  \quad+\quad
  \text{tag-subcode selection}\\
  \quad+\quad
  \text{residue arithmetic for transversal phases}.
\end{gathered}
  \label{eq:three-layers}
\end{equation}
The first two layers are GCQC structure.  The third layer is the SS--LP or transversal-gate
constraint.  The latter guarantees phase growth but says nothing by itself about distance.

The dyadic Bell lift is the smallest nonadditive instance of this template.  The two tag subcodes
are one-dimensional:
\begin{equation}
  T_0=\Span\{\ketphi\},
  \qquad
  T_1=\Span\{\ketpsi\}.
  \label{eq:bell-tag-subcodes}
\end{equation}
The carry bit selects $T_0$ or $T_1$, with a logical-label-dependent parity shift for the top
sector.  The tag is very small, so it cannot be invisible to all weight-two errors.  The detailed
Bell theorem compensates with the two-anchor induction.

\subsection{CSS linear tags and stabilizer-compatible lifts}

There is a stabilizer branch of the same idea.  Suppose a CSS code is described by classical
binary data $C\subseteq G\subseteq\F_2^n$, so that a logical coset state has the schematic form
\begin{equation}
  \ket{g+C}:=\frac{1}{\sqrt{|C|}}\sum_{c\in C}\ket{g+c}.
  \label{eq:css-coset-state}
\end{equation}
A linear carry-tag lift replaces the nonlinear carry map by a linear extension of the underlying
classical code.  The tag is then itself a CSS or stabilizer coset state, and the lifted support
remains a linear or affine-linear code.  The distance certificate is no longer a Bell-anchor KL
induction; it is a normalizer-distance calculation.  In a standard CSS notation the relevant
quantities are
\begin{equation}
  d_X=\min_{v\in G\setminus C}\wt(v),
  \qquad
  d_Z=\min_{z\in C^\perp\setminus G^\perp}\wt(z).
  \label{eq:css-distances}
\end{equation}
A stabilizer-compatible lift preserves distance at least $d$ when the lifted classical data make
both distances at least $d$.

This CSS branch is conceptually close to stabilizer divisible-code and Reed--Muller-type towers:
large linear tags preserve additivity and make the distance certificate classical, while small
nonlinear tags can leave the stabilizer category and require a direct KL proof.

\subsection{Relation to CSS diagonal-gate synthesis and surgery}

The Hu--Liang--Calderbank diagonal-gate synthesis framework can be interpreted as a Heisenberg- or
stabilizer-side version of the same separation between phase growth and distance repair.  The
framework combines three operations: concatenation, removal of $Z$-stabilizers, and addition of
$X$-stabilizers.  Concatenation changes the physical representation of the diagonal gate and can
raise the relevant level.  Removing $Z$-stabilizers splits or enlarges the logical sector so that
the desired induced diagonal action appears.  Adding $X$-stabilizers repairs distance losses.

In carry-tag language, concatenation introduces tag degrees of freedom, the stabilizer removal
step changes the sector decomposition, and the stabilizer addition step supplies a distance
certificate.  The distinction is that the present nonadditive Bell lift works directly at the
state-vector and KL level rather than in a stabilizer normalizer calculation.

\subsection{What GCQC does and does not guarantee}

The GCQC viewpoint explains the architecture but does not prove the main theorem by itself.  GCQC
organizes the code as sector data plus tag subcodes.  It also suggests general distance bounds when
inner and outer distances are known.  However, the dyadic Bell proof uses information finer than a
standard minimum-distance product bound: the exact Bell Pauli table, the unique anchor weights, the
absence of small strings in the opposite residue block, and a real-part cancellation between the
two logical sectors.  These are arithmetic and amplitude-level facts tied to the SS--LP residue
structure.

The useful lesson is methodological.  For nonadditive small-overhead lifts, one searches for compact
tag subcodes and proves a KL induction tailored to their low-weight error table.  For
stabilizer-preserving lifts, one searches for linear tag extensions and proves classical CSS
distance bounds.  Both are residue-compatible GCQC searches; they differ in the tag category and
in the distance certificate.

% ============================================================
\section{Comparison with related work}
\label{sec:comparison}
% ============================================================

\subsection{Stabilizer divisibility and Clifford-hierarchy towers}

Triorthogonal CSS codes and Bravyi--Haah magic-state-distillation codes show how divisibility
conditions on classical binary matrices can yield transversal non-Clifford gates in stabilizer
codes~\cite{BravyiHaah2012}.  Haah's generalized divisible-code towers extend this perspective: increasing the diagonal
gate level is possible while preserving distance, but the block length grows multiplicatively with
the target level~\cite{Haah2018}.  Rengaswamy--Calderbank--Newman--Pfister and Hu--Liang--Calderbank give a more
systematic algebraic description of when diagonal physical gates preserve stabilizer or CSS code
spaces and what logical gates they induce~\cite{RengaswamyCalderbankNewmanPfister2020,HuLiangCalderbank2022Channels,HuLiangCalderbank2022Divisible,HuLiangCalderbank2021Climbing}.

The dyadic Bell lift has a different profile.  It is generally nonadditive, fixes the target
distance at $3$, and increases the diagonal order by adding only two qubits per level.  This avoids
the stabilizer divisibility bottleneck by leaving the stabilizer category.  It does not contradict
stabilizer restrictions such as the Anderson--Jochym-O'Connor classification of diagonal
transversal gates for qubit stabilizer codes~\cite{AndersonJochymOConnor2016}, because those restrictions apply to stabilizer codes,
whereas the Bell carry map is generally nonlinear and non-stabilizer.

\subsection{Scaling comparison at fixed distance}
\label{sec:scaling-comparison}

The clearest quantitative distinction is obtained by fixing the target code dimension and
distance, and then asking how the block length grows as the diagonal rotation becomes finer.  In
this subsection the target is one encoded qubit, equivalently $K=2$, and distance at least $3$.
For a power-of-two projective order $2^\nu$ we write
\[
  D_{2^\nu}=\diag(1,e^{-2\pi i/2^\nu}).
\]
With this convention the $T$ gate, up to phase and adjoint convention, corresponds to $\nu=3$;
the $\sqrt T$ rotation corresponds to $\nu=4$; and $T^{1/4}$ corresponds to $\nu=5$.

The dyadic Bell lift is linear in this level parameter.  If a seed of length $n_0$ realizes
projective order $2^{\nu_0}$ and satisfies the hypotheses of
Theorem~\ref{thm:two-anchor-bell}, then after $r$ Bell lifts
\begin{equation}
  m_r=2^r m_0=2^{\nu_0+r},
  \qquad
  n_r=n_0+2r.
  \label{eq:dyadic-scaling-r}
\end{equation}
Equivalently, for level $\nu\ge \nu_0$,
\begin{equation}
  n(\nu)=n_0+2(\nu-\nu_0).
  \label{eq:dyadic-scaling-level}
\end{equation}
Thus each phase-doubling step costs exactly two additional physical qubits.

This scaling should be compared with stabilizer-code mechanisms.  The punctured quantum
Reed--Muller sequence gives distance-$3$ stabilizer codes with transversal diagonal gates at
successively higher levels; with the convention above, the corresponding lengths are
$15,31,63,\ldots$, i.e. $2^{\nu+1}-1$ up to the standard indexing convention.  Haah's generalized
divisible-code towers are more flexible in distance and input data, but the level-raising step has
multiplicative overhead; at fixed $d=3$ this gives schematic growth of order $3^\nu$
in the level parameter~\cite{Haah2018}.  Nonadditive permutation-invariant/spin-code constructions
give small individual codes, for example a distance-$3$ transversal-$T$ construction at
$n=11$, but they do not provide a dyadic ``check two rungs, get all rungs'' recursion
with constant two-qubit overhead~\cite{KubischtaTeixeira2023Exotic,KubischtaTeixeira2024PI}.

\begin{table}[t]
\centering
\small
\renewcommand{\arraystretch}{1.18}
\begin{tabular}{lcccc}
\hline\hline
\textbf{target level} & \textbf{Reed--Muller} & \textbf{Haah tower} & \textbf{PI/spin} & \textbf{dyadic Bell} \\
\hline
$\nu=3$ \;($T$)        & $15$  & $15$       & $11$ & $7$ \\
$\nu=4$ \;($\sqrt T$) & $31$  & $\sim45$  & no tower & $9$ \\
$\nu=5$ \;($T^{1/4}$) & $63$  & $\sim135$ & no tower & $11$ \\
$\nu=6$ \;($T^{1/8}$) & $127$ & $\sim405$ & no tower & $13$ \\
\hline
general $\nu$ & $2^{\nu+1}-1$ & $O(3^\nu)$ at $d=3$ & individual searches & $n_0+2(\nu-\nu_0)$ \\
\hline\hline
\end{tabular}
\caption{Representative block-length scaling at fixed $K=2$ and distance at least $3$.  The
Reed--Muller and Haah columns are stabilizer-code constructions.  The PI/spin column records known
small nonadditive constructions but no level-by-level recursive tower.  The displayed dyadic
numbers use a seed at $(n_0,\nu_0)=(7,3)$; the formula in the last column applies to any seed
satisfying Theorem~\ref{thm:two-anchor-bell}.}
\label{tab:scaling-comparison}
\end{table}

The structural reason for the different scaling is that stabilizer divisibility methods must
increase the divisibility order of large linear supports.  In contrast, the Bell lift records only
the new dyadic carry bit in a two-qubit tag.  Residue refinement proves the phase doubling, while
the two-anchor Knill--Laflamme induction proves that the weight-$\le2$ error-detection conditions
continue to hold after the finite base check.  Level growth is therefore paid for additively rather
than multiplicatively.

This comparison is deliberately narrow.  The dyadic Bell theorem fixes the distance target at
$3$ and keeps one encoded qubit, so its asymptotic rate tends to zero.  Stabilizer and CSS
constructions can target larger distance, and recent algebraic-geometry constructions obtain
asymptotically good quantum codes with transversal non-Clifford gates
\cite{GolowichGuruswami2024,Nguyen2024}.  The dyadic ladder is therefore best understood as a
small-block-length, fixed-distance, nonadditive mechanism for growing diagonal transversal order,
not as a replacement for high-rate or growing-distance stabilizer-code theory.

The same additive-overhead statement also applies outside the pure power-of-two setting.  Starting
from a certified seed of projective order $m_0$, the $r$-th dyadic Bell lift has order $2^r m_0$
and length $n_0+2r$.  If $m_0$ has an odd factor, the resulting diagonal rotations are not indexed
only by the usual power-of-two Clifford-hierarchy levels, but the carry-tag arithmetic and the
distance-propagation theorem are unchanged.

\subsection{Permutation-invariant and spin-code constructions}

Permutation-invariant and spin-code constructions provide another route to nonadditive codes with
large transversal gate groups~\cite{KubischtaTeixeira2023Exotic,KubischtaTeixeira2024PI}.  Their organizing principle is representation-theoretic: the code
space is chosen inside a symmetric or spin representation with a prescribed finite group action.
The dyadic Bell lift instead uses subset-sum residues and carry tags.  The two approaches share the
central use of nonadditivity to move beyond stabilizer Clifford-hierarchy constraints, but they
organize the search space differently.  The Bell theorem supplies a recursive finite-verification
mechanism: once the seed and first lift are checked, the entire ladder follows from the two-anchor
induction.

\subsection{SS--LP and recent nonadditive diagonal-transversal searches}

The SS--LP framework introduced in recent nonadditive transversal-gate searches fixes diagonal
residue sectors first and then solves amplitude constraints~\cite{ZhangWuHuangZeng2025}.  Carry-tag lifts add a recursive
operation to that framework.  Rather than searching independently at each target modulus, one can
start from a seed satisfying the two-anchor hypotheses and generate a whole dyadic ladder.  The
logical phase growth is guaranteed by residue refinement, and the distance is guaranteed by the
Bell theorem's finite KL certificate.

The comparison with automated or data-driven SS--LP searches is therefore complementary.  Search is
well suited to finding seeds and discovering new residue patterns.  The lift theorem turns suitable
seeds into infinite families and explains why the same distance certificate continues to hold after
the first lift.

\subsection{GCQC perspective}

Generalized concatenated quantum codes are broad enough to include both stabilizer and nonadditive
families.  The carry-tag template should be viewed as GCQC with an additional arithmetic constraint
coming from transversal diagonal gates.  Relative to standard GCQC, the new feature is that the
outer-sector labels are subset-sum carries and the inner/tag sectors are chosen to refine diagonal
phases.  Relative to stabilizer concatenation, the new feature is that the tag can be nonlinear and
very small, with distance proved by a direct KL induction rather than by CSS normalizer distances.

% ============================================================
\section{Summary}
\label{sec:summary}
% ============================================================

The SS--LP viewpoint separates transversal phase design from error correction.  Subset-sum residue
blocks force a tensor-product diagonal gate to act as a desired logical diagonal.  The KL equations
then decide whether the chosen supports and amplitudes form a code of the desired distance.

A carry-tag lift refines residue classes modulo $m$ to residue classes modulo $2m$.  The tag stores
the dyadic carry bit so that the lifted logical states land in the correct refined residue sectors.
This proves diagonal order growth, but not distance.  Distance requires a separate certificate.

The nonadditive dyadic Bell lift is a minimal two-qubit carry-tag construction.  Its distance-$3$
certificate is the two-anchor Bell theorem: strong subset-sum arithmetic, unique bottom and top
anchors, seed-plus-first-lift KL checks, Bell-pair matrix elements, small-weight support separation,
and anchor-leakage cancellation together imply an infinite distance-$3$ ladder.  The construction
is naturally understood as a residue-compatible GCQC, with the Bell pair as a small nonadditive tag
and the theorem as the distance certificate.

\begin{thebibliography}{99}

\bibitem{KnillLaflamme1997}
E.~Knill and R.~Laflamme,
\emph{Theory of quantum error-correcting codes},
Physical Review A \textbf{55}, 900--911 (1997).

\bibitem{CalderbankShor1996}
A.~R. Calderbank and P.~W. Shor,
\emph{Good quantum error-correcting codes exist},
Physical Review A \textbf{54}, 1098--1105 (1996).

\bibitem{Steane1996}
A.~M. Steane,
\emph{Error correcting codes in quantum theory},
Physical Review Letters \textbf{77}, 793--797 (1996).

\bibitem{EastinKnill2009}
B.~Eastin and E.~Knill,
\emph{Restrictions on transversal encoded quantum gate sets},
Physical Review Letters \textbf{102}, 110502 (2009).

\bibitem{AndersonJochymOConnor2016}
J.~T. Anderson and T.~Jochym-O'Connor,
\emph{Classification of transversal gates in qubit stabilizer codes},
Quantum Information \& Computation \textbf{16}, 771--802 (2016); arXiv:1409.8320.

\bibitem{BravyiHaah2012}
S.~Bravyi and J.~Haah,
\emph{Magic-state distillation with low overhead},
Physical Review A \textbf{86}, 052329 (2012).

\bibitem{Haah2018}
J.~Haah,
\emph{Towers of generalized divisible quantum codes},
Physical Review A \textbf{97}, 042327 (2018); arXiv:1709.08658.

\bibitem{RengaswamyCalderbankNewmanPfister2020}
N.~Rengaswamy, R.~Calderbank, M.~Newman, and H.~D. Pfister,
\emph{On optimality of CSS codes for transversal $T$},
arXiv:1910.09333.

\bibitem{HuLiangCalderbank2022Channels}
J.~Hu, Q.~Liang, and R.~Calderbank,
\emph{Designing the quantum channels induced by diagonal gates},
Quantum \textbf{6}, 802 (2022); arXiv:2109.13481.

\bibitem{HuLiangCalderbank2022Divisible}
J.~Hu, Q.~Liang, and R.~Calderbank,
\emph{Divisible codes for quantum computation},
arXiv:2204.13176.

\bibitem{HuLiangCalderbank2021Climbing}
J.~Hu, Q.~Liang, and R.~Calderbank,
\emph{Climbing the diagonal Clifford hierarchy},
arXiv:2110.11923.

\bibitem{GrasslShorSmithSmolinZeng2009}
M.~Grassl, P.~W. Shor, G.~Smith, J.~A. Smolin, and B.~Zeng,
\emph{Generalized concatenated quantum codes},
arXiv:0901.1319.

\bibitem{KubischtaTeixeira2023Exotic}
E.~Kubischta and I.~Teixeira,
\emph{Family of quantum codes with exotic transversal gates},
Physical Review Letters \textbf{131}, 240601 (2023).

\bibitem{KubischtaTeixeira2024PI}
E.~Kubischta and I.~Teixeira,
\emph{Permutation-invariant quantum codes with transversal generalized phase gates},
arXiv:2310.17652.

\bibitem{ZhangWuHuangZeng2025}
C.~Zhang, Z.~Wu, S.~Huang, and B.~Zeng,
\emph{Transversal gates in nonadditive quantum codes},
arXiv:2504.20847.

\bibitem{GolowichGuruswami2024}
L.~Golowich and V.~Guruswami,
\emph{Asymptotically good quantum codes with transversal non-Clifford gates},
arXiv:2408.09254.

\bibitem{Nguyen2024}
Q.~T. Nguyen,
\emph{Good binary quantum codes with transversal CCZ gate},
arXiv:2408.10140.

\bibitem{HeLuZeng2025}
X.~He, S.~Lu, and B.~Zeng,
\emph{Co-designing quantum codes with transversal diagonal gates via multi-agent systems},
arXiv:2510.20728.

\end{thebibliography}

\end{document}
